\documentclass[10pt,a4paper,twocolumn]{article}

\usepackage[a4paper, top=1.2cm, bottom=1.2cm, left=1cm, right=1cm, columnsep=0.8cm]{geometry}
\usepackage[T1]{fontenc}
\usepackage[expansion=false]{microtype}
\usepackage{amsmath, amssymb, amsfonts, braket}
\usepackage{graphicx}
\usepackage{tikz-cd}
\usepackage{booktabs}
\usepackage{array}
\usepackage{enumitem}
\usepackage[dvipsnames,table]{xcolor}
\usepackage[most]{tcolorbox}
\usepackage{titlesec}

\setlist{itemsep=2pt, topsep=3pt, leftmargin=1.2em, parsep=0pt}
\setlength{\parskip}{4pt}
\setlength{\parindent}{0pt}
\linespread{1.05}

% ── Farben ──────────────────────────────────────────────
\definecolor{c1}{RGB}{25,95,155}     % Vektorräume – blau
\definecolor{c2}{RGB}{110,45,150}    % Lin. Unabh. – lila
\definecolor{c3}{RGB}{15,120,70}     % Basis & Dim – grün
\definecolor{c4}{RGB}{175,80,0}      % Norm – orange
\definecolor{c5}{RGB}{155,25,55}     % Fourier – rot
\definecolor{c6}{RGB}{0,110,125}     % Abbildungen – teal
\definecolor{c7}{RGB}{222,93,131}     % Kern u. Bild – blush
\definecolor{c8}{RGB}{181,166,66}     % Isomorphismen – brass
\definecolor{c9}{RGB}{189,51,164}     % Basiswechsel – byzantine
\definecolor{c10}{RGB}{0,191,255}     % Eigenvektor u. Eigenwerte – capri
\definecolor{c11}{RGB}{113,188,120}     % Diagonaliseren v. Matrizen – fern
\definecolor{cn}{RGB}{80,80,80}      % Notes – grau
\definecolor{detail}{RGB}{100,100,100}
\definecolor{defbg}{RGB}{232,242,255}
\definecolor{tipbg}{RGB}{255,247,220}
\definecolor{recbg}{RGB}{230,247,236}
\definecolor{detbg}{RGB}{248,248,248}

\newcommand{\seccol}{black}
\newcommand{\setseccol}[1]{\renewcommand{\seccol}{#1}}

% ── Section-Format ──────────────────────────────────────
\newcommand{\makesecheader}[1]{%
	\noindent\colorbox{\seccol}{\makebox[\columnwidth][l]{%
			\color{white}\bfseries\small\strut\enspace\thesection\enspace\textbar\enspace #1\enspace%
	}}%
}
\titleformat{\section}[block]{\normalfont}{}{0pt}{\makesecheader}
\titlespacing*{\section}{0pt}{8pt}{4pt}

\titleformat{\subsection}[block]{\normalfont\bfseries\small}{}{0pt}{\color{\seccol}\rule[-.3ex]{2.5pt}{9pt}\enspace}
\titlespacing*{\subsection}{0pt}{6pt}{2pt}

\titleformat{\subsubsection}[block]{\normalfont\bfseries\footnotesize\color{detail}}{}{0pt}{}
\titlespacing*{\subsubsection}{0pt}{4pt}{1pt}

% ── Box-Stile ───────────────────────────────────────────
\tcbset{
	base/.style={boxrule=0.4pt, arc=2.5pt, left=5pt, right=5pt, top=4pt, bottom=4pt,
		fonttitle=\bfseries\small, enhanced},
	defbox/.style={base, colback=defbg, colframe=c1!60,
		attach boxed title to top left, boxed title style={colback=c1!60}, coltitle=white},
	tipbox/.style={base, colback=tipbg, colframe=c4!70, notitle},
	recbox/.style={base, colback=recbg, colframe=c3!70,
		coltitle=white},
	detbox/.style={base, colback=detbg, colframe=gray!30, notitle,
		fontupper=\small\color{detail}},
}

% Kleiner Detail-Text
\newcommand{\detail}[1]{{\small\color{detail}#1}}

\begin{document}
	\twocolumn[{%
		\centering\vspace{-2pt}
		{\large\bfseries Zusammenfassung — Lineare Algebra 2}
		\enspace{\small\color{gray}Woche 1–7}
		\vspace{3pt}\hrule\vspace{5pt}
	}]
	
	%═══════════════════════════════════════════════════
	\setseccol{c1}
	\section{Vektorräume \& Unterräume}
	
	\begin{tcolorbox}[defbox, title={Def. 5.1.1 — Vektorraum}]
		$\mathbb{K}$-Vektorraum: Menge $V$ mit $+$ und $\cdot$, 8 Axiome:
		\begin{itemize}
			\item \textbf{Addition (i–iv):} Assoziativität, Nullvektor, Negatives, Kommutativität
			\item \textbf{Skalarmult. (v–vi):} $(\lambda\mu)v=\lambda(\mu v)$,\ $1\cdot v=v$
			\item \textbf{Distributiv (vii–viii):} $\lambda(v+w)=\lambda v+\lambda w$
		\end{itemize}
	\end{tcolorbox}
	
	\detail{Beispiele: $\mathbb{R}^n$, $\mathbb{C}^n$, $\mathbb{K}^{m\times n}$, $P_n$ (Polynome), $C([a,b])$ (stetige Fkt.)}
	
	\subsection{Unterräume}
	
	\begin{tcolorbox}[defbox, title={Def. 5.2.1 — Unterraum}]
		$U\subseteq V$ ist Unterraum $\iff$
		\begin{itemize}[leftmargin=1em]
			\item[i)] $u,v\in U \Rightarrow u+v\in U$
			\item[ii)] $\lambda\in\mathbb{K},\,u\in U \Rightarrow \lambda u\in U$
		\end{itemize}
	\end{tcolorbox}
	
	\begin{tcolorbox}[tipbox]
		\textbf{GO-Kriterium:} $0\notin U \Rightarrow$ sofort kein Unterraum!
	\end{tcolorbox}
	
	\detail{%
		\textbf{Unterräume:} Geraden/Ebenen durch Ursprung, $\ker(A)$, symmetr.\ Matrizen.\\
		\textbf{Kein Unterraum:} $Ax=c$ ($c\neq0$), $\det(A)=1$, invertierbare Matrizen.%
	}
	
	%═══════════════════════════════════════════════════
	\setseccol{c2}
	\section{Lineare Unabhängigkeit}
	
	\begin{tcolorbox}[defbox, title={Definition}]
		$v_1,\ldots,v_n$ \textbf{lin.\ unabhängig} $\iff$
		$$\sum_{i=1}^n\lambda_i v_i = 0 \implies \lambda_i=0\ \forall i$$
		Sonst \textbf{linear abhängig}.
	\end{tcolorbox}
	
	\begin{tcolorbox}[recbox, title={Rezept via Gauss}]
		Vektoren als Spalten $\to$ Gauss $\to$ Rang:\\
		Rang $=$ Anz.\ Vektoren $\Rightarrow$ \textbf{unabhängig}\\
		Rang $<$ Anz.\ Vektoren $\Rightarrow$ \textbf{abhängig}
	\end{tcolorbox}
	
	\detail{Abhängig wenn: Nullvektor enthalten, kollineare Vektoren, oder ein Vektor ist Linearkombination der anderen. In $\mathbb{R}^2$ max.\ 2, in $\mathbb{R}^3$ max.\ 3 lin.\ unabh.\ Vektoren.}
	
	\begin{tcolorbox}[detbox]
		\textbf{Bsp.\ unabhängig:}\ $v_1=\!\binom{1}{0}$, $v_2=\!\binom{0}{1}$: Rang $=2=$ Anz.\ $\Rightarrow$ unabh.\\
		\textbf{Bsp.\ abhängig:}\ $v_1=\!\binom{1}{2}$, $v_2=\!\binom{2}{4}$: Rang $=1<2$ $\Rightarrow$ abh.
	\end{tcolorbox}
	
	\begin{tcolorbox}[detbox]
		\textbf{Lin. Unabhängigkeit von Matrizen:} Jede Matrix $A_i \in \mathbb{R}^{m\times n}$ als Spaltenvektor schreiben (Spalten untereinander, $mn$ Einträge). Dann wie Vektoren behandeln: als Spalten in eine Matrix $\to$ Gauss $\to$ Rang bestimmen.\\
		Oder direkt: $\lambda_1 A_1 + \cdots + \lambda_k A_k = 0$ komponentenweise lösen (jeder Eintrag gibt eine Gleichung).
	\end{tcolorbox}
	
	\subsection{Lineare Hülle \& Erzeugendensystem}
	
	$$\operatorname{span}(v_1,\ldots,v_n):=\!\left\{v\in V\;\middle|\;v=\sum_i\lambda_i v_i\right\}$$
	
	\detail{$S$ heisst \textbf{Erzeugendensystem} von $V$ falls $\operatorname{span}(S)=V$. Eine lin.\ Hülle muss nicht lin.\ unabhängig sein — nur als Basis.}
	
	%═══════════════════════════════════════════════════
	\setseccol{c3}
	\section{Basis und Dimension}
	
	\begin{tcolorbox}[defbox, title={Def. 5.4.1 — Basis}]
		$B\subset V$ ist Basis $\iff$ lin.\ unabhängig \textbf{und} $\operatorname{span}(B)=V$.\\
		Kleinstmögliches Erzeugendensystem.
	\end{tcolorbox}
	
	\begin{tcolorbox}[defbox, title={Def. 5.4.7 — Dimension}]
		$\dim_\mathbb{K}(V)$ = Anzahl Vektoren in einer Basis (eindeutig).\\
		$\dim(\mathbb{R}^n)=n$,\quad $\dim(P_2)=3$,\quad $\dim(\{0\})=0$
	\end{tcolorbox}
	
	\begin{tcolorbox}[tipbox]
		\textbf{Koordinatenvektor (Satz 5.4.9):}
		$v=\sum x_i b_i \;\leftrightarrow\; x=(x_1,\ldots,x_n)^\top$ \textbf{eindeutig}.\\
		Koordinaten hängen von der Basiswahl ab!
	\end{tcolorbox}
	
	\detail{Standardbasis $\mathbb{R}^n$: $e_1,\ldots,e_n$ (Einheitsvektoren). Bsp.\ $P_2$: $\{1,x,x^2\}$ oder $\{3,x^2,-2x\}$ — viele Basen möglich.\\
		\textbf{Basisergänzungssatz:} Lin.\ unabh.\ Vektoren $\to$ zur Basis ergänzen; Erzeugendensystem $\to$ Basis herausstreichen.}
	
	%═══════════════════════════════════════════════════
	\setseccol{c4}
	\section{Norm \& Skalarprodukt}
	
	\subsection{Vektornormen}
	
	\begin{tcolorbox}[defbox, title={Normen}]
		$\|v\|_p=\bigl(\sum|v_i|^p\bigr)^{1/p}$\quad (p-Norm)\\
		$\|v\|_2=\sqrt{\sum v_i^2}$\quad (Euklidisch)\quad $\|v\|_\infty=\max_i|v_i|$
	\end{tcolorbox}
	
	\detail{Bsp.: $\|\binom{3}{4}\|_2=5$;\quad $\|(3,-7,2)^\top\|_\infty=7$}
	
	\subsection{Matrixnormen}
	
	\begin{tcolorbox}[defbox, title={Matrixnormen}]
		Frobenius: $\|A\|_F=\sqrt{\sum_{ij}|a_{ij}|^2}$\\
		Zeilensumme: $\|A\|_\infty=\max_i\sum_j|a_{ij}|$\\
		Spaltensumme: $\|A\|_1=\max_j\sum_i|a_{ij}|$
	\end{tcolorbox}
	
	\detail{Für $A=\bigl(\begin{smallmatrix}1&-2\\3&4\end{smallmatrix}\bigr)$: $\|A\|_F=\sqrt{30}$,\ $\|A\|_\infty=7$,\ $\|A\|_1=6$}
	
	\subsection{Funktionsnorm}
	\detail{$L^p$-Norm: $\|f\|_{L^p}=\bigl(\int_\mathbb{D}|f|^p dx\bigr)^{1/p}$. \quad Bsp.: $\|x\|_{L^2[0,1]}=\frac{1}{\sqrt{3}}$}
	
	\subsection{Skalarprodukt}
	
	\begin{tcolorbox}[defbox, title={Skalarprodukte}]
		$\mathbb{R}^n$: $\langle x,y\rangle=\sum x_iy_i$\quad
		$\mathbb{C}^n$: $\langle w,z\rangle=\sum\overline{w_i}z_i$\\
		Matrizen: $\langle A,B\rangle=\operatorname{tr}(A^TB)$\\
		$L^2$: $\langle f,g\rangle=\int_\mathbb{D}\overline{f(x)}g(x)\,dx$
	\end{tcolorbox}
	
	\begin{tcolorbox}[tipbox]
		\textbf{Induzierte Norm:} $\|v\|=\sqrt{\langle v,v\rangle}$.\quad Jedes Skalarprodukt $\Rightarrow$ Norm (nicht umgekehrt).
	\end{tcolorbox}
	
	\detail{Bsp.: $\langle(1,2),(3,4)\rangle=11$;\quad $\langle f,g\rangle_{L^2}$ für $f=x,g=x^2$ auf $[0,1]$: $\int_0^1 x^3dx=\frac{1}{4}$}
	
	\subsection{Orthonormalbasen (ONB)}
	
	\begin{tcolorbox}[defbox, title={Def. 5.6.8 — ONB}]
		$\{e_1,\ldots,e_n\}$ ist ONB $\iff \langle e_i,e_j\rangle=\delta_{ij}$\\
		(paarweise orthogonal + normiert)
	\end{tcolorbox}
	
	\begin{tcolorbox}[tipbox]
		\textbf{Vorteil:} Koordinaten ohne LGS: $x_i=\langle e_i,v\rangle$\\
		$v=\sum_{i=1}^n\langle e_i,v\rangle\,e_i$
	\end{tcolorbox}
	
	\subsection{Gram-Schmidt}
	
	\begin{tcolorbox}[recbox, title={Algorithmus}]
		\textbf{1.\ Orthogonalisieren:}
		$$v_1=w_1,\quad v_i=w_i-\sum_{k=1}^{i-1}\frac{\langle v_k,w_i\rangle}{\langle v_k,v_k\rangle}v_k$$
		\textbf{2.\ Normieren:} $e_i=v_i/\|v_i\|$
	\end{tcolorbox}
	
	%═══════════════════════════════════════════════════
	\setseccol{c5}
	\section{Fourier-Reihen}
	
	\detail{Idee: Fourier-Koeffizienten = Koordinaten in ONB des $L^2$-Raums. $L^2([0,T])$ mit $\langle f,g\rangle=\frac{1}{T}\int_0^T\overline{f}g\,dt$.}
	
	\begin{tcolorbox}[tipbox]
		\textbf{$\omega$ aus Signal ablesen:} $\omega=\operatorname{ggT}$ aller Frequenzen.
		$\omega=\pi\cdot\frac{\operatorname{ggT(Zähler)}}{\operatorname{kgV(Nenner)}}$\\
		\detail{Bsp.: $\frac{\pi}{3},\frac{5\pi}{6},\frac{8\pi}{3}\Rightarrow\omega=\frac{\pi}{6}$, also $n=2,5,16$.}
	\end{tcolorbox}
	
	\subsection{Reelle Schreibweise}
	
	\begin{tcolorbox}[defbox, title={Formel}]
		$$f(t)=\frac{a_0}{2}+\sum_{n=1}^\infty\bigl(a_n\cos(n\omega t)+b_n\sin(n\omega t)\bigr)$$
	\end{tcolorbox}
	
	\textbf{Koeffizienten:}
	$$a_0=\frac{2}{T}\int_0^T\!f\,dt, \quad a_n=\frac{2}{T}\int_0^T\!f\cos(n\omega t)\,dt$$
	$$b_n=\frac{2}{T}\int_0^T\!f\sin(n\omega t)\,dt$$
	
	\begin{tcolorbox}[tipbox]
		$f$ \textbf{gerade} $\Rightarrow b_n=0$.\quad $f$ \textbf{ungerade} $\Rightarrow a_n=0$.\\
		$a_0/2=0$: kein DC-Anteil (Mittelwert null).
	\end{tcolorbox}
	
	\subsection{Komplexe Schreibweise}
	
	\begin{tcolorbox}[defbox, title={Formel}]
		$f(t)=\sum_{k=-\infty}^{\infty}c_k e^{ik\omega t},\quad c_k=\frac{1}{T}\int_0^T\!f(t)e^{-ik\omega t}dt$
	\end{tcolorbox}
	
	\begin{tcolorbox}[tipbox]
		\textbf{DC-Anteil (komplex):} $c_0 = \frac{1}{T}\int_0^T f(t)\,dt = \frac{a_0}{2}$
	\end{tcolorbox}
	
	\detail{Basisfunktionen $e^{ik\omega t}$ bilden vollst.\ ONS. Approximation: $f_N(t)=\sum_{k=-N}^N c_k e^{ik\omega t}$.}
	
	\subsection{Amplituden-Phasen-Form}
	
	\begin{tcolorbox}[defbox, title={Formel}]
		$f(t)=\frac{a_0}{2}+\sum_{n=1}^\infty A_n\cos(n\omega t+\varphi_n)$\\
		$A_n=\sqrt{a_n^2+b_n^2}$,\quad $\varphi_n=\arg(c_n)$,\quad $|c_n|=\frac{A_n}{2}$
	\end{tcolorbox}
	
	\subsection{Umrechnung zwischen Schreibweisen}
	
	\begin{tcolorbox}[recbox, title={Umrechnungsformeln}]
		\textbf{Ampl.-Phase $\to$ Reell:}\ $a_k=A_k\cos\varphi_k$,\ $b_k=-A_k\sin\varphi_k$\\
		\textbf{Reell $\to$ Komplex:}\ $c_k=\frac{a_k}{2}-\frac{b_k}{2}i$,\ $c_{-k}=\overline{c_k}$\\
		\textbf{Komplex $\to$ Reell:}\ $a_k=2\operatorname{Re}(c_k)$,\ $b_k=-2\operatorname{Im}(c_k)$
	\end{tcolorbox}
	
	\detail{Alle drei Schreibweisen = dasselbe Signal. $c_{-k}$ nur in komplexer Schreibweise, immer $=\overline{c_k}$.}
	
	\subsection{Konvergenz}
	
	\begin{itemize}
		\item \textbf{$L^2$:} Immer — Fläche zw.\ $f_N$ und $f\to 0$.
		\item \textbf{Dirichlet:} Endlich viele Unstetigkeitsstellen $\Rightarrow$ punktweise.
		\item \textbf{Gleichmässig:} $f$ stetig + stückw.\ stetig diffbar.
		\item \textbf{Gibbs:} Überschwingungen an Sprüngen bleiben (auch grosse $N$).
	\end{itemize}
	
	%═══════════════════════════════════════════════════
	\setseccol{c6}
	\section{Lineare Abbildungen}
	
	\begin{tcolorbox}[tipbox]
		\textbf{Spalten von $A$ = Bilder der Basisvektoren!}\\
		Jede $m\times n$-Matrix $\leftrightarrow$ lineare Abbildung $\mathbb{K}^n\to\mathbb{K}^m$.
	\end{tcolorbox}
	
	\begin{tcolorbox}[defbox, title={Def. 6.1.1 — Lineare Abbildung}]
		$f:V\to W$ linear $\iff$
		\begin{itemize}[leftmargin=1em]
			\item[i)] $f(v_1+v_2)=f(v_1)+f(v_2)$
			\item[ii)] $f(\lambda v)=\lambda f(v)$
		\end{itemize}
		\textbf{Folgerung:} $f(0)=0$ — wenn $f(0)\neq0$: nicht linear!
	\end{tcolorbox}
	
	\detail{Beispiele: Matrizen, Ableitung $\frac{d}{dx}:P_2\to P_1$, Projektion, Spiegelung, Drehung.\\
		Nicht linear: $f(x)=Ax+c$ mit $c\neq0$.}
	
	\subsubsection{Spezielle Matrizen — Erkennung}
	
	\begin{tcolorbox}[defbox, title={Übersicht}]
		\begin{tabular}{@{}lccc@{}}
			\toprule
			& $\det$ & $A^TA{=}I$? & $A^2{=}A$? \\
			\midrule
			Drehung & $+1$ & ja & nein \\
			Spiegelung & $-1$ & ja & nein \\
			Projektion & $0$ & nein & ja \\
			Streckung $\alpha$ & $\alpha^2$ & nein & nein \\
			\bottomrule
		\end{tabular}
	\end{tcolorbox}
	
	\textbf{Drehung} um $\varphi$ im $\mathbb{R}^2$:
	$$O=\begin{pmatrix}\cos\varphi&-\sin\varphi\\\sin\varphi&\cos\varphi\end{pmatrix}\in SO(2),\quad O^TO=I,\quad\det O=1$$
	
	\textbf{Spiegelung} an $x_2$-Achse:
	$$S=\begin{pmatrix}-1&0\\0&1\end{pmatrix},\quad\det S=-1$$
	
	\textbf{Projektion} auf $x_1$-$x_2$-Ebene:
	$$P=\begin{pmatrix}1&0&0\\0&1&0\end{pmatrix},\quad\det P=0,\quad P^2=P$$
	
	\textbf{Streckung} um Faktor $\alpha$:
	$$S=\begin{pmatrix}\alpha&0\\0&\alpha\end{pmatrix},\quad\det S=\alpha^2$$
	
	\detail{Trick: $\cos(x+\tfrac{\pi}{2})=-\sin x$,\quad $\cos(x-\tfrac{\pi}{2})=+\sin x$}
	
	\subsection{Matrix einer linearen Abbildung}
	
	\begin{tcolorbox}[recbox, title={Rezept: Matrix aufstellen — Normalfall}]
		\begin{enumerate}[leftmargin=1.5em, itemsep=2pt]
			\item Basen $B_V=\{v_1,\ldots,v_n\}$, $B_W=\{w_1,\ldots,w_m\}$ wählen
			\item Bilder $f(v_j)$ der Basisvektoren \textbf{einzeln berechnen}
			\item $f(v_j)$ als Linearkombination in Basis $B_W$ ausdrücken
			\item Koeffizienten $\to$ $j$-te \textbf{Spalte} von $A$
		\end{enumerate}
	\end{tcolorbox}
	
	\begin{tcolorbox}[tipbox]
		\textbf{Sonderfall: Bilder der Basisvektoren nicht direkt bekannt!}\\
		Wenn man $f(e_1), f(e_2)$ nicht kennt, aber $f$ auf anderen Vektoren weiss:\\[3pt]
		\textbf{1.} Diese Vektoren als Linearkombination der Basis schreiben:\\
		\quad $\binom{1}{2} = e_1 + 2e_2 \;\Rightarrow\; f\binom{1}{2} = f(e_1) + 2f(e_2) = \binom{7}{0}$\\[3pt]
		\textbf{2.} Gleiches für zweiten bekannten Vektor:\\
		\quad $\binom{-2}{2} = -2e_1 + 2e_2 \;\Rightarrow\; -2f(e_1) + 2f(e_2) = \binom{4}{-6}$\\[3pt]
		\textbf{3.} LGS lösen (Gauss oder Addition/Subtraktion der Gleichungen) um $f(e_1)$ und $f(e_2)$ zu finden.\\[3pt]
		\textbf{4.} Bilder als Spalten in die Matrix schreiben.
	\end{tcolorbox}
	
	\detail{Bsp.: $f(e_1)=\binom{1}{2}$, $f(e_2)=\binom{3}{-1}$ $\Rightarrow$ $A=\bigl(\begin{smallmatrix}1&3\\2&-1\end{smallmatrix}\bigr)$.\\
		\textbf{Merke:} Die Matrix hängt von der Basiswahl ab. Koordinatenformel: $y=Ax$.}
	
	\begin{center}
		\begin{tikzcd}[column sep=1.8em, row sep=1.2em]
			V \arrow[r, "f"] \arrow[d, "B_V"'] & W \arrow[d, "B_W"] \\
			\mathbb{K}^n \arrow[r, "A"'] & \mathbb{K}^m
		\end{tikzcd}
	\end{center}
	
	\detail{Bsp.\ $\frac{d}{dx}:P_2\to P_1$, Standardbasen: $(1)'=0$, $(x)'=1$, $(x^2)'=2x$ $\Rightarrow$ $A=\bigl(\begin{smallmatrix}0&1&0\\0&0&2\end{smallmatrix}\bigr)$.\\
		\textbf{Koordinatenformel:} $y=Ax$.\quad \textbf{Eindeutigkeit:} Zu jeder Matrix $A$ genau eine lin.\ Abbildung (bzgl.\ Basen) und umgekehrt. Matrix hängt von Basiswahl ab!}
	
	\subsection{Verkettung \& Matrixprodukt}
	
	\begin{tcolorbox}[defbox, title={Satz 6.3.2}]
		Matrix von $(g\circ f)$ = $B\cdot A$\\
		(erst $f$ mit Matrix $A$, dann $g$ mit Matrix $B$)
	\end{tcolorbox}
	
	\begin{tcolorbox}[tipbox]
		$AB\neq BA$ im Allgemeinen — \textbf{nicht kommutativ}!\\
		Assoziativität gilt: $(C\circ B)\circ A=C\circ(B\circ A)$.
	\end{tcolorbox}
	
	\setseccol{c7}
	\section{Kern und Bild einer Lineare Abbildung}
	
	\begin{tcolorbox}[defbox, title={Def. 7.1.1 — Vektorraum}]
		$\mathbb{K}$-Vektorräume: $V$ und $W$ und $f : V \to W$:
		\begin{itemize}
			\item $\ker(f) := \{v \in V | f(v)=0\}$
			\item $\operatorname{im}(f) := \{w \in W | w = f(v) \text{ für ein } v \in V\}$
		\end{itemize}
	\end{tcolorbox}
	
	\detail{%
		\textbf{Kern:} Menge der Vektoren, die von einer linearen
		Abbildung auf den Nullvektor abgebildet werden.\\
		\textbf{Bild:} Menge der Vektoren,
		auf die die Vektoren durch die lineare Abbildung abgebildet werden.\\
	}
	
	\begin{tcolorbox}[defbox, title={Def. 7.1.2 — Vektorraum}]
		\begin{itemize}
			\item $\ker(f)$ ist ein Unterraum von $V$
			\item $\operatorname{im}(f)$ ist ein Unterraum von $W$
		\end{itemize}
	\end{tcolorbox}
	
	\begin{tcolorbox}[recbox, title={Rezept: Berechnung von ker() und im():}]
		$\mathbb{K}$-Vektorräume: $V$ und $W$ und $f : V \to W$.\\
		Matrix $A$ von $f$ bezüglich gewählter Basen $B_V = \{v_1 , . . . , v_n\}$ und $BW = \{w_1 , . . . , w_m\}$ gegeben.\\
		\\
		\begin{tikzcd}[column sep=1.8em, row sep=1.2em]
			V \arrow[r, "f"] \arrow[d, "B_V"'] & W \arrow[d, "B_W"] \\
			\mathbb{K}^n \arrow[r, "A"'] & \mathbb{K}^m
		\end{tikzcd}
		\\
		
		\textbf{1.\ Kern:}
		\begin{enumerate}[leftmargin=1.5em, itemsep=2pt]
			\item \textbf{Löse das homogene lineare Gleichungssystem Ax = 0 (Gauss)}
			\item Nullvektor \textbf{immer} im Kern.
			\item Lösungsmenge von $Ax =0$ immer ein $(n - r)$-
			dimensionaler Unterraum von $K^n$
			\item r = rang(A) die Anzahl Pivotelemente im Endschema von $Ax = 0$
			\item maximaler Rang $\to r = n$: Kern trivial und enthält nur den Nullvektor.
		\end{enumerate}
		\textbf{2.\ Bild:}
		\begin{enumerate}[leftmargin=1.5em, itemsep=2pt]
			\item Spalten von $A$ spannen das Bild von $A$ auf.
			\item linear \textbf{abhängigen} Spalten wegstreichen.
			\item Spalten die übrigbleiben, bilden eine Basis von $\operatorname{im}(f)$.
			\item Gauss-Algorithmus: Spalten mit Pivot $\to$ Basis des Bildes
			von $f$.
			\item $\dim_{\mathbb{K}} (\operatorname{im}(f)) = \operatorname{rang}(A)$
		\end{enumerate}
	\end{tcolorbox}
	
	\begin{tcolorbox}[defbox, title={Def. 7.1.3 — Rangsatz}]
		$f : V \to W$ eine lineare Abbildung, wobei $dim_{\mathbb{K}} (V ) = n < \infty$. Dann gilt:
		$$\dim_{\mathbb{K}} (\ker(f )) + \dim_{\mathbb{K}} (\operatorname{im}(f)) = \dim_{\mathbb{K}} (V ).$$
	\end{tcolorbox}
	
	\setseccol{c8}
	\section{Isomorphismen (Inversion)}
	
	\setseccol{c9}
	\section{Basiswechsel}
	
	\setseccol{c10}
	\section{Eigenvektor und Eigenwerte}
	
	\setseccol{c11}
	\section{Diagonalisieren von Matrizen}
	
	%═══════════════════════════════════════════════════
	\setseccol{cn}
	\section{Notes}
	
	\detail{%
		\textbf{Senkrechter Vektor:} $(a,b)^\perp=(-b,a)$ — Werte tauschen, ein Minus.\\
		\textbf{Determinante \& Fläche:} $|\det(A)|$ = Fläche des Parallelogramms der Spaltenvektoren.\\
	}
	
	\detail{%
		\textbf{Matrix:} m = Zeilebn ; n = Spalten ; $m \times n$
	}
	
	\detail{%
		\textbf{Rezept Abbildungsmatrix:} (1) Basisvektoren von $B_V$ abbilden, (2) Resultate in Basis $B_W$ ausdrücken.\\
		Jede $m\times n$-Matrix erzeugt lineare Abbildung $\mathbb{R}^n\to\mathbb{R}^m$.%
	}
	
	
	
\end{document}